\section{Analysis of Model-\texorpdfstring{$B_2^1$}{B21} (\texorpdfstring{$\gamma_1>0$}{gamma1>0} head-of-line, \texorpdfstring{$\gamma_2=\theta_1=\theta_2=0$}{gamma2=theta1=theta2=0})}
\label{sec:model_b21}

Model-$B_2^1$ is a deliberately simplified variant of the one-way jockeying model
Model-$B_2$. In Model-$B_2$ \emph{every} waiting class-$1$ customer may jockey to the
class-$2$ queue, so the aggregate jockeying rate out of a state grows linearly with the
queue length, $\gamma_1 n_1$. Here, by contrast, only the \emph{customer at the head of
Queue~1} --- the one that would be the next class-$1$ customer to enter service --- is
allowed to jockey, at the fixed rate $\gamma_1$ whenever $n_1\ge 1$. The aggregate
jockeying rate is therefore $\gamma_1\,\mathbf{1}_{\{n_1\ge 1\}}$ rather than
$\gamma_1 n_1$.

This single change is decisive for tractability. In Model-$B_2$ the linear coefficient
$\gamma_1 n_1$ produces, upon forming the joint PGF, a term proportional to
$xy(x-y)\,\partial P/\partial x$, turning the fundamental equation into a first-order
linear ODE in $x$ that requires the integrating-factor machinery of
Section~\ref{sec:model_b2}. The flat rate $\gamma_1\,\mathbf{1}_{\{n_1\ge1\}}$, on the
other hand, contributes only the \emph{algebraic} factor $\gamma_1 y(x-y)$ to the
coefficient of $P(x,y)$; \emph{no derivative of $P$ appears}. The fundamental equation is
then purely algebraic and can be solved by the Kernel Method exactly as in Model-$A$
(Section~\ref{sec:model_a}). The analysis below mirrors the analytical treatment of
Model-$A$ step by step, and the resulting PGF turns out to be \emph{structurally
identical} to that of Model-$A$, the only difference being a shift in the characteristic
roots.

\textbf{Stability.} Jockeying relocates a waiting customer between the two queues but does
not remove anyone from the system; the total number of customers is therefore conserved by
the jockeying mechanism, exactly as in Model-$A$. Since there is no abandonment, the chain
is positive recurrent if and only if the total traffic intensity satisfies
\begin{equation}
    \rho = \frac{\lambda_1+\lambda_2}{\mu} < 1,
    \label{eq:B21:stability}
\end{equation}
with $\rho_1=\lambda_1/\mu$ governing the priority class.

\subsection{Balance equations}

Because the jockeying rate is now constant in $n_1$ rather than linear, the balance
equations are \emph{not} obtained by zeroing parameters in the general
Lemma~\ref{lem:gen:fundamental_equation} (whose jockeying coefficient is $\gamma_1 n_1$);
they must be stated directly. The jockeying transition $(n_1,n_2)\to(n_1-1,n_2+1)$ occurs
at the fixed rate $\gamma_1$ whenever $n_1\ge1$. On the state space $S$ the global balance
equations read
\begin{align}
    \begin{split}
        &\left[\lambda_1+\lambda_2+\mu+\gamma_1\right]\pi(n_1,n_2) \\
        &\quad=\quad \mu\,\pi(n_1+1,n_2)
          +\lambda_1\pi(n_1-1,n_2)
          +\lambda_2\pi(n_1,n_2-1)
          +\gamma_1\pi(n_1+1,n_2-1)
    \end{split}
    &&n_1\ge1,\;n_2\ge1, \label{eq:B21:interior} \\
    \begin{split}
        &\left[\lambda_1+\lambda_2+\mu+\gamma_1\right]\pi(n_1,0)
          \;=\; \mu\,\pi(n_1+1,0)
          +\lambda_1\pi(n_1-1,0)
    \end{split}
    &&n_1\ge1,\;n_2=0, \label{eq:B21:xboundary} \\
    \begin{split}
        &\left[\lambda_1+\lambda_2+\mu\right]\pi(0,n_2) \\
        &\quad=\quad \mu\,\pi(1,n_2)
          +\gamma_1\pi(1,n_2-1)
          +\lambda_2\pi(0,n_2-1)
          +\mu\,\pi(0,n_2+1)
    \end{split}
    &&n_1=0,\;n_2\ge1, \label{eq:B21:yboundary} \\
    \begin{split}
        &\left[\lambda_1+\lambda_2+\mu\right]\pi(0,0)
          = (\lambda_1+\lambda_2)\pi_0
          + \mu\,\pi(1,0)
          + \mu\,\pi(0,1),
    \end{split} &&\label{eq:B21:zero} \\
    \begin{split}
        &(\lambda_1+\lambda_2)\pi_0 = \mu\,\pi(0,0).
    \end{split} &&\label{eq:B21:idle}
\end{align}
Two structural remarks are in order. First, jockeying is an \emph{out}-transition from any
state with $n_1\ge1$ (hence the $+\gamma_1$ in the left-hand sides of
\eqref{eq:B21:interior} and~\eqref{eq:B21:xboundary}) and an \emph{in}-transition into any
state with $n_2\ge1$ from the state $(n_1+1,n_2-1)$ (hence the $\gamma_1\pi(n_1+1,n_2-1)$
in~\eqref{eq:B21:interior} and the $\gamma_1\pi(1,n_2-1)$ in~\eqref{eq:B21:yboundary}).
Second, when $n_2=0$ there is no jockeying in-flow, since it would have to originate from
the non-existent state $(n_1+1,-1)$; and when $n_1=0$ there is no jockeying out-flow. The
remaining transitions --- Poisson arrivals at $\lambda_1,\lambda_2$ and service
completions at $\mu$ --- are exactly those of Model-$A$.

\subsection{Fundamental PGF equation}

\begin{lemma}[Fundamental equation for Model-$B_2^1$]\label{lem:B21:fundamental}
Assume positive recurrence, so that the stationary distribution exists and
$P(x,y)=\sum_{n_1,n_2\ge0}\pi(n_1,n_2)x^{n_1}y^{n_2}$ is analytic on the closed unit
polydisc. Then $P(x,y)$ and the boundary function $P_y(y)=P(0,y)$ satisfy
\begin{equation}
    \begin{split}
        &\left[(\lambda_1+\lambda_2+\mu)xy - \mu y - \lambda_1 x^2y - \lambda_2 xy^2
               + \gamma_1 y(x-y)\right]P(x,y) \\
        &\quad=\quad \left[\mu(x-y) + \gamma_1 y(x-y)\right]P_y(y)
               - \mu x(1-y)\,\pi(0,0).
    \end{split}
    \label{eq:B21:fundamental}
\end{equation}
\end{lemma}

\begin{proof}
Apart from the jockeying terms, equations~\eqref{eq:B21:interior}--\eqref{eq:B21:idle}
coincide with the balance equations of Model-$A$. Multiplying each balance equation by
$x^{n_1}y^{n_2}$, summing over its domain and multiplying through by $xy$, the
Model-$A$ part reproduces exactly the left- and right-hand sides of the Model-$A$
fundamental equation~\eqref{eq:A:fundamental},
\begin{equation*}
    \left[ (\lambda_1+\lambda_2+\mu)xy - \mu y - \lambda_1 x^2y - \lambda_2 xy^2 \right] P(x,y)
    = \mu(x-y) P_y(y) - \mu x(1-y)\,\pi(0,0),
\end{equation*}
so it remains only to collect the jockeying contribution. Its \emph{out}-flow appears in
every state with $n_1\ge1$:
\begin{equation*}
    \gamma_1\sum_{n_1\ge1}\sum_{n_2\ge0}\pi(n_1,n_2)\,x^{n_1}y^{n_2}
    = \gamma_1\bigl[P(x,y)-P_y(y)\bigr],
\end{equation*}
while its \emph{in}-flow enters every state with $n_2\ge1$ from $(n_1+1,n_2-1)$:
\begin{equation*}
    \gamma_1\sum_{n_1\ge0}\sum_{n_2\ge1}\pi(n_1+1,n_2-1)\,x^{n_1}y^{n_2}
    = \gamma_1\,\frac{y}{x}\sum_{m_1\ge1}\sum_{m_2\ge0}\pi(m_1,m_2)\,x^{m_1}y^{m_2}
    = \gamma_1\,\frac{y}{x}\bigl[P(x,y)-P_y(y)\bigr],
\end{equation*}
where we substituted $m_1=n_1+1,\ m_2=n_2-1$. The net jockeying contribution (out minus
in) is therefore $\gamma_1\bigl(1-\tfrac{y}{x}\bigr)\bigl[P(x,y)-P_y(y)\bigr]$; multiplying
by $xy$ gives $\gamma_1 y(x-y)\bigl[P(x,y)-P_y(y)\bigr]$. Adding this to the Model-$A$
equation and moving the $P_y$ part to the right-hand side yields~\eqref{eq:B21:fundamental}.
\end{proof}

\noindent\textit{Remark.}
Equation~\eqref{eq:B21:fundamental} contains \emph{no derivative} of $P$. The flat
jockeying rate has contributed the algebraic factor $\gamma_1 y(x-y)$ to the coefficient of
$P(x,y)$ where Model-$B_2$, with its linear rate $\gamma_1 n_1$, contributes the
convective term $\gamma_1 xy(x-y)\,\partial P/\partial x$. Setting $\gamma_1=0$
recovers~\eqref{eq:A:fundamental} verbatim.

\subsection{Closed-form solution}

The following theorem is the main result of this section; its proof follows the
analytical (Kernel Method) treatment of Model-$A$ in~\S\ref{sssec:A:analytical}.

\begin{theorem}[Closed-form PGF for Model-$B_2^1$]\label{thm:B21:PGF}
Consider the two-class non-preemptive priority queue with Poisson arrival rates
$\lambda_1,\lambda_2$, common exponential service rate $\mu$, and head-of-line class-$1$
jockeying at rate $\gamma_1>0$ (with $\gamma_2=\theta_1=\theta_2=0$). Under the stability
condition~\eqref{eq:B21:stability}, the joint PGF $P(x,y)=\mathbb{E}[x^{N_1}y^{N_2}]$ is
the rational function
\begin{equation}
    P(x,y) \;=\; \frac{\mu\,\rho(1-\rho)\,(1-y)}
    {\lambda_1\,\bigl(x^*(y)-y\bigr)\,\bigl(x^+(y)-x\bigr)},
    \label{eq:B21:PGF}
\end{equation}
where $x^*(y)$ and $x^+(y)$ are the two roots of the kernel quadratic
\begin{equation}
    \lambda_1 x^2 - \bigl[\lambda_1+\lambda_2(1-y)+\mu+\gamma_1\bigr]x + (\mu+\gamma_1 y) = 0,
    \label{eq:B21:quad}
\end{equation}
with $x^*(y)$ the root inside the unit disk. This is the \emph{same rational form} as
Model-$A$ (Theorem~\ref{thm:A:PGF}); the two differ only in the kernel, where $\mu$ has been
replaced by $\mu+\gamma_1 y$.
\end{theorem}

\begin{corollary}\label{cor:B21:pi0}
Under~\eqref{eq:B21:stability}, the empty-state probabilities of Model-$B_2^1$ coincide with
those of Model-$A$,
\begin{equation}
    \pi_0 = 1-\rho, \qquad \pi(0,0) = \rho(1-\rho).
    \label{eq:B21:pi0_pi00}
\end{equation}
In particular $P(1,1)=1-\pi_0=\rho$ and the diagonal generating function satisfies
$P(z,z)=\pi(0,0)/(1-\rho z)$.
\end{corollary}

\begin{proof}[Proof of Corollary~\ref{cor:B21:pi0}]
Set $x=y=z$ in the fundamental equation~\eqref{eq:B21:fundamental}. The jockeying factor
$\gamma_1 y(x-y)$ vanishes on the diagonal --- a direct consequence of the fact that
jockeying conserves the total number of customers --- and so does the term
$\mu(x-y)P_y(y)$, leaving
\begin{equation*}
    \bigl[(\lambda_1+\lambda_2+\mu)z^2 - \mu z - (\lambda_1+\lambda_2)z^3\bigr]P(z,z)
    = -\mu z(1-z)\,\pi(0,0).
\end{equation*}
Dividing by $z$ and factoring the bracket as $-(\lambda z-\mu)(z-1)$, with
$\lambda=\lambda_1+\lambda_2$,
\begin{equation*}
    (\lambda z-\mu)(1-z)\,P(z,z) = -\mu(1-z)\,\pi(0,0).
\end{equation*}
Cancelling $(1-z)\ne0$ for $z\in[0,1)$ gives the diagonal identity
\begin{equation}
    P(z,z) = \frac{\mu\,\pi(0,0)}{\mu-\lambda z} = \frac{\pi(0,0)}{1-\rho z}.
    \label{eq:B21:Pzz}
\end{equation}
Setting $z=1$ and using the normalisation $P(1,1)=1-\pi_0$ yields
$1-\pi_0=\pi(0,0)/(1-\rho)$. The idle balance~\eqref{eq:B21:idle} gives
$\lambda\pi_0=\mu\pi(0,0)$, i.e.\ $\pi(0,0)=\rho\pi_0$; substituting,
$(1-\pi_0)(1-\rho)=\rho\pi_0$, hence $\pi_0=1-\rho$ and $\pi(0,0)=\rho(1-\rho)$.
\end{proof}

\subsubsection{Analytical approach for Model-\texorpdfstring{$B_2^1$}{B21}}\label{sssec:B21:analytical}

\begin{proof}[Proof of Theorem~\ref{thm:B21:PGF}]
We apply the \emph{Kernel Method} to the fundamental equation~\eqref{eq:B21:fundamental}.
Fix $y\in(0,1]$ and seek the values of $x$ at which the coefficient of $P(x,y)$ --- the
\emph{kernel} --- vanishes. Since $P(x,y)$ is a PGF it is bounded on the closed unit disk,
so at any such root the right-hand side of~\eqref{eq:B21:fundamental} must vanish as well.
Dividing the kernel by $y\ne0$,
\begin{equation*}
    0 = (\lambda_1+\lambda_2+\mu)x - \mu - \lambda_1 x^2 - \lambda_2 xy + \gamma_1(x-y),
\end{equation*}
and rearranging into a quadratic in $x$ gives exactly~\eqref{eq:B21:quad},
\begin{equation*}
    \lambda_1 x^2 - \bigl[\lambda_1+\lambda_2(1-y)+\mu+\gamma_1\bigr]x + (\mu+\gamma_1 y) = 0,
\end{equation*}
whose two roots are
\begin{equation}
    x^{\pm}(y) = \frac{A(y) \pm \sqrt{A(y)^2 - 4\lambda_1(\mu+\gamma_1 y)}}{2\lambda_1},
    \qquad A(y):=\lambda_1+\lambda_2(1-y)+\mu+\gamma_1.
    \label{eq:B21:xstar}
\end{equation}
This is precisely the Model-$A$ kernel quadratic~\eqref{eq:A:xstar} with $\mu$ replaced by
$\mu+\gamma_1 y$. We label the root carrying the negative sign $x^*(y)$ and the other
$x^+(y)$. By Vieta's formulas the product of the roots is
$x^*(y)\,x^+(y)=(\mu+\gamma_1 y)/\lambda_1$.

To locate the roots relative to the unit disk we evaluate the kernel quadratic, written as
$f(x):=\lambda_1 x^2 - A(y)x + (\mu+\gamma_1 y)$, at $x=1$:
\begin{equation*}
    f(1) = \lambda_1 - \bigl[\lambda_1+\lambda_2(1-y)+\mu+\gamma_1\bigr] + (\mu+\gamma_1 y)
         = -(1-y)\,(\lambda_2+\gamma_1) < 0
    \qquad\text{for } y\in[0,1).
\end{equation*}
Since $f$ is an upward parabola ($\lambda_1>0$) with $f(1)<0$, the point $x=1$ lies strictly
between its two roots, so $x^*(y)<1<x^+(y)$ for all $y\in[0,1)$. Hence $x^*(y)$ is the unique
root inside the unit disk, and at $y=1$ one checks $f(1)=0$ with $x^*(1)=1$ and
$x^+(1)=(\mu+\gamma_1)/\lambda_1>1$ (the latter exceeding $1$ because
$\lambda_1<\mu<\mu+\gamma_1$ under stability). For later use, we record
$\tfrac{d}{dy}x^*(1)$. The cleanest route is implicit differentiation of
$f(x,y):=\lambda_1 x^2-A(y)x+(\mu+\gamma_1 y)=0$ at the point $(x,y)=(1,1)$, where
$f_y=\lambda_2 x+\gamma_1$ and $f_x=2\lambda_1 x-A(y)$ evaluate to
$f_y(1,1)=\lambda_2+\gamma_1$ and $f_x(1,1)=\lambda_1-\mu-\gamma_1$, so
$\tfrac{d}{dy}x^*(1)=-f_y/f_x$:
\begin{equation}
    \frac{d}{dy}x^*(1) = \frac{\lambda_2+\gamma_1}{\mu+\gamma_1-\lambda_1}.
    \label{eq:B21:xstar_deriv}
\end{equation}
Unlike Model-$A$'s~\eqref{eq:A:xstar_deriv}, the numerator carries an extra $\gamma_1$:
the constant term $\mu+\gamma_1 y$ of the kernel~\eqref{eq:B21:quad} is itself
$y$-dependent here, whereas in Model-$A$ (and in Model-$C_2^1$ below) the constant term is
independent of $y$.

Setting $x=x^*(y)$ in~\eqref{eq:B21:fundamental} and equating the right-hand side to zero,
\begin{equation*}
    \bigl[\mu(x^*(y)-y) + \gamma_1 y(x^*(y)-y)\bigr]P_y(y) - \mu x^*(y)(1-y)\,\pi(0,0) = 0,
\end{equation*}
that is, $(x^*(y)-y)(\mu+\gamma_1 y)P_y(y) = \mu x^*(y)(1-y)\pi(0,0)$, whence the boundary
function is
\begin{equation}
    P_y(y) = \frac{\mu\,x^*(y)\,(1-y)\,\pi(0,0)}
                  {(\mu+\gamma_1 y)\,\bigl(x^*(y)-y\bigr)}.
    \label{eq:B21:Py}
\end{equation}
As a consistency check, set $y=0$ in~\eqref{eq:B21:Py}: the factor $\mu+\gamma_1 y$ becomes
$\mu$ and $1-y$ becomes $1$, so
$P_y(0)=\mu\,x^*(0)\,\pi(0,0)/\bigl(\mu\,x^*(0)\bigr)=\pi(0,0)$ (using $x^*(0)>0$, the
smaller of the two positive roots), recovering the boundary value $P_y(0)=\pi(0,0)$.

Finally, substitute~\eqref{eq:B21:Py} and $\pi(0,0)=\rho(1-\rho)$
(Corollary~\ref{cor:B21:pi0}) back into~\eqref{eq:B21:fundamental}. Writing the kernel as
$-\lambda_1 y\,(x-x^*(y))(x-x^+(y))$ and simplifying the right-hand side,
\begin{align*}
    \bigl[\mu(x-y)+\gamma_1 y(x-y)\bigr]P_y(y) - \mu x(1-y)\,\pi(0,0)
    &= \mu(1-y)\,\pi(0,0)\left[\frac{(x-y)\,x^*(y)}{x^*(y)-y}-x\right] \\
    &= \mu(1-y)\,\pi(0,0)\,\frac{y\,\bigl(x-x^*(y)\bigr)}{x^*(y)-y},
\end{align*}
where the bracket collapses because
$(x-y)x^*-x(x^*-y)=y(x-x^*)$. Dividing by the factored kernel, the common factor
$\bigl(x-x^*(y)\bigr)$ cancels and we obtain
\begin{equation*}
    P(x,y) = \frac{\mu(1-y)\,\pi(0,0)\,y\,(x-x^*(y))/(x^*(y)-y)}
                  {-\lambda_1 y\,(x-x^*(y))(x-x^+(y))}
           = \frac{\mu\,\rho(1-\rho)\,(1-y)}
                  {\lambda_1\,\bigl(x^*(y)-y\bigr)\,\bigl(x^+(y)-x\bigr)},
\end{equation*}
which is~\eqref{eq:B21:PGF} and concludes the proof.
\end{proof}

\subsubsection{Probabilistic approach for Model-\texorpdfstring{$B_2^1$}{B21}}\label{sssec:B21:probabilistic}

For Model-$A$, the boundary function $P_y(y)$ admits an independent derivation by the
\emph{Pollaczek--Khinchine} (PK) formula (\S\ref{sssec:A:probabilistic}): conditional on
$N_1=0$, class~$2$ experiences an $M/G/1$ queue whose service times are i.i.d.\ copies of a
class-$1$ busy period, and the PK formula then delivers $P_y(y)$. \emph{This route does not
transfer to Model-$B_2^1$.} Head-of-line jockeying injects class-$1$ customers into Queue~2
at the instants when $N_1\ge1$; the resulting arrival stream into Queue~2 is the
superposition of the exogenous Poisson stream $\lambda_2$ and a state-dependent jockeying
stream. This merged stream is \emph{neither} Poisson \emph{nor} independent of the system
state --- it is correlated with $N_1$. Since the PK formula requires PASTA, which in turn
demands both a Poisson arrival process \emph{and} the lack-of-anticipation property, and the
merged input violates both, class~$2$ does not see a renewal $M/G/1$ input and no PK
representation of $P_y(y)$ is available. The analytical derivation
of~\S\ref{sssec:B21:analytical} is consequently the only route, and we retain this
subsection solely to record why the probabilistic one is unavailable --- in the same spirit
as the discussion of Models~$B$ and~$B_2$.

\subsection{Limits and sanity checks}

As $\gamma_1\to0^+$ the kernel quadratic~\eqref{eq:B21:quad} reduces to Model-$A$'s
quadratic~\eqref{eq:A:xstar}, so $x^*(y)\to x_A^*(y)$, and formula~\eqref{eq:B21:PGF}
reduces to the factored form of Theorem~\ref{thm:A:PGF}; the empty-state probabilities
already equal those of Model-$A$ by Corollary~\ref{cor:B21:pi0}. Setting $x=y=1$
in~\eqref{eq:B21:PGF} through the diagonal identity~\eqref{eq:B21:Pzz} gives
$P(1,1)=\rho=1-\pi_0$, confirming normalisation, and the diagonal
$P(z,z)=\pi(0,0)/(1-\rho z)$ is exactly that of the combined $M/M/1(\lambda,\mu)$
queue --- as it must be, since jockeying conserves the total customer count.

\subsection{Mean queue lengths and mean waiting times}

\subsubsection{Class-1 marginal PGF and mean queue length}

Setting $y=1$ in the fundamental equation~\eqref{eq:B21:fundamental}, the term
$\mu x(1-y)\pi(0,0)$ vanishes and the coefficient of $P(x,1)$ factors as
$-(x-1)\bigl(\lambda_1 x-(\mu+\gamma_1)\bigr)$, while the right-hand side becomes
$(x-1)(\mu+\gamma_1)P_y(1)$. Cancelling $(x-1)$ for $x\ne1$,
\begin{equation}
    P(x,1) \;=\; \frac{(\mu+\gamma_1)\,P_y(1)}{\mu+\gamma_1-\lambda_1 x}
            \;=\; \frac{P_y(1)}{1-\rho_B\,x},
    \qquad \rho_B := \frac{\lambda_1}{\mu+\gamma_1}.
    \label{eq:B21:Px1}
\end{equation}
This is a geometric class-$1$ marginal: the priority queue behaves as an $M/M/1$ queue whose
effective \emph{clearance rate} is $\mu+\gamma_1$, since a waiting class-$1$ customer at the
head of Queue~1 departs that queue either by entering service (rate $\mu$) or by jockeying
(rate $\gamma_1$). The constant $P_y(1)$ is pinned by the normalisation
$P(1,1)=1-\pi_0=\rho$, giving $P_y(1)=\rho(1-\rho_B)$, so that
$P(x,1)=\rho(1-\rho_B)/(1-\rho_B x)$. Differentiating by the quotient rule and evaluating at
$x=1$,
\begin{equation}
    \mathbb{E}[N_1] \;=\; \left.\frac{d}{dx}P(x,1)\right|_{x=1}
    \;=\; \frac{\rho(1-\rho_B)\,\rho_B}{(1-\rho_B)^2}
    \;=\; \frac{\rho\,\rho_B}{1-\rho_B}.
    \label{eq:B21:EN1}
\end{equation}
This is Model-$A$'s formula~\eqref{eq:A:EN1} with the priority load $\rho_1$ replaced by the
\emph{jockeying-reduced} load $\rho_B=\lambda_1/(\mu+\gamma_1)\le\rho_1$.

\subsubsection{Total queue length from the diagonal}

Because jockeying conserves the total number of customers, the combined process
$N_1+N_2$ (together with the in-service customer) is the queue-length process of an ordinary
$M/M/1(\lambda,\mu)$ queue, and the diagonal generating
function~\eqref{eq:B21:Pzz} is $P(z,z)=\rho(1-\rho)/(1-\rho z)$. Since
$\tfrac{d}{dz}P(z,z)\big|_{z=1}=\mathbb{E}[N_1]+\mathbb{E}[N_2]$, differentiating gives
\begin{equation}
    \mathbb{E}[N_1]+\mathbb{E}[N_2]
    \;=\; \left.\frac{d}{dz}\,\frac{\rho(1-\rho)}{1-\rho z}\right|_{z=1}
    \;=\; \frac{\rho^2(1-\rho)}{(1-\rho)^2}
    \;=\; \frac{\rho^2}{1-\rho},
    \label{eq:B21:EN}
\end{equation}
identical to the combined $M/M/1$ result of Model-$A$~\eqref{eq:A:EN2}.

\subsubsection{Class-2 mean queue length by subtraction}

Subtracting~\eqref{eq:B21:EN1} from~\eqref{eq:B21:EN},
\begin{equation}
    \begin{split}
        \mathbb{E}[N_2]
        &\;=\; \frac{\rho^2}{1-\rho} - \frac{\rho\,\rho_B}{1-\rho_B}
        \;=\; \frac{\rho\bigl[\rho(1-\rho_B)-\rho_B(1-\rho)\bigr]}{(1-\rho)(1-\rho_B)} \\[4pt]
        &\;=\; \frac{\rho\,(\rho-\rho_B)}{(1-\rho)(1-\rho_B)}.
    \end{split}
    \label{eq:B21:EN2}
\end{equation}
The excess load admits a transparent decomposition,
\begin{equation*}
    \rho-\rho_B
    = \frac{\lambda_1+\lambda_2}{\mu} - \frac{\lambda_1}{\mu+\gamma_1}
    = \frac{(\lambda_1+\lambda_2)(\mu+\gamma_1)-\lambda_1\mu}{\mu(\mu+\gamma_1)}
    = \frac{\lambda_2\mu+\lambda_1\gamma_1+\lambda_2\gamma_1}{\mu(\mu+\gamma_1)}
    = \rho_2 + \frac{\lambda_1\gamma_1}{\mu(\mu+\gamma_1)},
\end{equation*}
so $\rho-\rho_B$ exceeds the bare class-$2$ load $\rho_2$ precisely by the load that
jockeying transfers from Queue~1 into Queue~2. At $\gamma_1=0$ it collapses to $\rho_2$
and~\eqref{eq:B21:EN2} recovers Model-$A$'s
$\mathbb{E}[N_2]=\rho\rho_2/[(1-\rho_1)(1-\rho)]$.

\subsubsection{Mean waiting times via Little's Law}

We apply \textbf{Little's Law} to each queue separately, recalling that $N_i$ counts only
the customers \emph{waiting}, so $\mathbb{E}[N_i]=\Lambda_i\,\mathbb{E}[W_i]$ with
$\Lambda_i$ the long-run rate at which customers \emph{enter} Queue~$i$ and $\mathbb{E}[W_i]$
their mean residence time in that queue (from entry until they leave it, by service entry or
by jockeying).

Every class-$1$ customer enters Queue~1 exactly once, by exogenous arrival, so
$\Lambda_1=\lambda_1$ and
\begin{equation}
    \mathbb{E}[W_1] \;=\; \frac{\mathbb{E}[N_1]}{\lambda_1}
    \;=\; \frac{\rho\,\rho_B}{\lambda_1(1-\rho_B)}
    \;=\; \frac{\rho}{\mu+\gamma_1-\lambda_1}.
    \label{eq:B21:EW1}
\end{equation}
Compared with Model-$A$'s $\mathbb{E}[W_1]=\rho/(\mu-\lambda_1)$~\eqref{eq:A:EW1}, the
denominator is enlarged by $\gamma_1$: head-of-line jockeying drains the priority queue
faster and shortens class-$1$ waits.

Queue~2 receives a \emph{mixed} input: the exogenous Poisson stream at rate $\lambda_2$
together with the jockeying stream. The latter fires at rate $\gamma_1$ whenever $N_1\ge1$,
so its long-run rate is $\gamma_1\,\mathbb{P}(N_1\ge1)$. The probability that a class-$1$
customer is waiting follows from the geometric marginal~\eqref{eq:B21:Px1} by subtracting
the $N_1=0$ mass from the busy mass,
\begin{equation*}
    \mathbb{P}(N_1\ge1) = P(1,1)-P(0,1) = \rho - \rho(1-\rho_B) = \rho\rho_B,
\end{equation*}
since $P(1,1)=1-\pi_0=\rho$ and $P(0,1)=P_y(1)=\rho(1-\rho_B)$. The total entrance rate into
Queue~2 is therefore $\Lambda_2=\lambda_2+\gamma_1\rho\rho_B$, and Little's Law gives the
mean residence time per Queue-2 entrant,
\begin{equation}
    \mathbb{E}[W_2] \;=\; \frac{\mathbb{E}[N_2]}{\Lambda_2}
    \;=\; \frac{\rho\,(\rho-\rho_B)}
               {(1-\rho)(1-\rho_B)\,\bigl(\lambda_2+\gamma_1\rho\rho_B\bigr)}.
    \label{eq:B21:EW2}
\end{equation}
As $\gamma_1\to0^+$ the jockeying stream disappears, $\Lambda_2\to\lambda_2$ and
$\rho_B\to\rho_1$, and~\eqref{eq:B21:EW2} recovers Model-$A$'s
$\mathbb{E}[W_2]=\rho/[\mu(1-\rho_1)(1-\rho)]$~\eqref{eq:A:EW2}.


% ═════════════════════════════════════════════════════════════════════════════
